\documentclass[12pt]{article}

\usepackage[utf8]{inputenc}
\usepackage[T1]{fontenc}
\usepackage{lmodern}
\usepackage[spanish]{babel}
\usepackage{booktabs}
\usepackage{amsmath}
\usepackage{forest}
\usepackage{float}
\usepackage{listings}
\usepackage{xcolor}
\usepackage{tikz}

\definecolor{codegreen}{rgb}{0,0.6,0}
\definecolor{codegray}{rgb}{0.5,0.5,0.5}
\definecolor{codepurple}{rgb}{0.58,0,0.82}
\definecolor{backcolour}{rgb}{0.95,0.95,0.92}

\lstdefinestyle{mystyle}{
    backgroundcolor=\color{backcolour},   
    commentstyle=\color{codegreen},
    keywordstyle=\color{magenta},
    numberstyle=\tiny\color{codegray},
    stringstyle=\color{codepurple},
    basicstyle=\ttfamily\footnotesize,
    breakatwhitespace=false,         
    breaklines=true,                 
    captionpos=b,                    
    keepspaces=true,                 
    numbers=left,                    
    numbersep=5pt,                  
    showspaces=false,                
    showstringspaces=false,
    showtabs=false,                  
    tabsize=2
}

\lstset{style=mystyle}

\sloppy
\setlength{\parindent}{0pt}

\begin{document}

% Título y materia
\begin{center}
  {\LARGE \textbf{Guía de Ejercicios\\Validación de Modelos}}\\[0.5em]
  {Analítica de Datos, Universidad de San Andrés}
\end{center}

Si encuentran algún error en el documento o hay alguna duda, mandenmé un mail a rodriguezf@udesa.edu.ar y lo revisamos.

\section{Ejercicios}

\subsection{Ejercicios Conceptuales}

\subsubsection{Verdadero o Falso}

\begin{enumerate}
    \item La accuracy es siempre la mejor métrica para evaluar un modelo de clasificación.
    \item Un modelo con AUC = 0.3 es peor que uno aleatorio.
    \item La precisión mide la proporción de verdaderos positivos entre todas las predicciones positivas.
    \item El recall y la especificidad suman siempre 1.
    \item En un problema médico es más importante maximizar el recall que la precisión.
    \item La curva ROC grafica precisión versus recall.
    \item Un F1-Score alto garantiza que tanto precisión como recall son altos.
    \item La matriz de confusión solo se puede usar para problemas binarios.
    \item AUC = 0.5 indica un clasificador perfecto.
    \item La curva PR es más informativa que ROC cuando las clases están desbalanceadas.
    \item Una tasa de falsos positivos alta es siempre peor que una baja.
    \item La especificidad es más importante que el recall en sistemas de justicia.
    \item El False Discovery Rate y la precisión suman 1.
    \item En problemas multiclase no existen falsos positivos.
    \item Un modelo puede tener alta accuracy pero baja utilidad práctica.
    \item El MSE penaliza más los errores grandes que el MAE.
    \item Un R² negativo indica que el modelo es peor que predecir la media.
    \item El RMSE tiene las mismas unidades que la variable objetivo.
    \item El MAPE puede causar problemas cuando hay valores cero en los datos reales.
    \item Un modelo con R² = 0.8 explica el 80\% de la varianza de la variable objetivo.
\end{enumerate}

\subsubsection{Verdadero o Falso (Cross Validation)}

\begin{enumerate}
    \item Cross validation permite obtener una estimación más robusta del rendimiento del modelo.
    \item En k-fold cross validation, el dataset se divide en k particiones exactamente iguales.
    \item Leave-one-out cross validation es siempre la mejor opción para evaluar modelos.
    \item Cross validation se puede usar tanto para seleccionar hiperparámetros como para evaluar el rendimiento final.
    \item En 5-fold cross validation, cada observación se usa exactamente 4 veces para entrenar y 1 vez para validar.
    \item Stratified cross validation es importante para datasets desbalanceados.
    \item Cross validation siempre reduce el overfitting del modelo.
    \item Un mayor número de folds siempre es mejor en cross validation.
\end{enumerate}

\subsubsection{Multiple Choice}

\begin{enumerate}
    \item ¿Cuál es la diferencia principal entre precisión y recall?
    \begin{enumerate}
        \item No hay diferencia, son sinónimos
        \item Precisión se enfoca en predicciones positivas, recall en casos positivos reales
        \item Recall es para problemas binarios, precisión para multiclase
        \item Precisión es más importante que recall
    \end{enumerate}
    
    \item En un modelo de detección de cáncer, ¿qué métrica priorizarías?
    \begin{enumerate}
        \item Accuracy
        \item Precisión
        \item Recall
        \item F1-Score
    \end{enumerate}
    
    \item ¿Qué representa un punto en la curva ROC?
    \begin{enumerate}
        \item Un modelo diferente
        \item Un umbral de clasificación diferente
        \item Una métrica diferente
        \item Un dataset diferente
    \end{enumerate}
    
    \item Si tengo un dataset con 95\% de clase negativa y 5\% de clase positiva, ¿qué métrica puede ser engañosa?
    \begin{enumerate}
        \item Accuracy
        \item Recall
        \item Precisión
        \item AUC
    \end{enumerate}
    
    \item ¿Cuándo usarías la curva PR en lugar de ROC?
    \begin{enumerate}
        \item Siempre
        \item Nunca
        \item Cuando las clases están balanceadas
        \item Cuando las clases están desbalanceadas
    \end{enumerate}
    
    \item ¿Qué significa AUC = 0.8?
    \begin{enumerate}
        \item 80\% de accuracy
        \item 80\% de probabilidad de ranquear correctamente un par positivo-negativo
        \item 80\% de precisión
        \item 80\% de recall
    \end{enumerate}
    
    \item ¿Cuál métrica es más sensible a valores atípicos?
    \begin{enumerate}
        \item MSE
        \item MAE
        \item RMSE
        \item R²
    \end{enumerate}
    
    \item ¿Qué significa un R² = 0.3?
    \begin{enumerate}
        \item El modelo explica el 30\% de la varianza
        \item El modelo tiene 30\% de accuracy
        \item El modelo es 30\% mejor que la media
        \item El modelo es 30\% peor que aleatorio
    \end{enumerate}
    
    \item ¿Cuál es la principal ventaja del MAPE?
    \begin{enumerate}
        \item Es más preciso que otras métricas
        \item Es fácil de interpretar como porcentaje
        \item No tiene problemas con valores cero
        \item Siempre es mejor que el MAE
    \end{enumerate}
\end{enumerate}

\subsubsection{Multiple Choice (Cross Validation)}

\begin{enumerate}
    \item ¿Cuál es la principal ventaja de cross validation sobre hold-out validation?
    \begin{enumerate}
        \item Es más rápido computacionalmente
        \item Usa toda la data para entrenamiento y validación
        \item No requiere dividir los datos
        \item Es más fácil de implementar
    \end{enumerate}
    
    \item En 10-fold cross validation, ¿qué porcentaje de los datos se usa para entrenar en cada fold?
    \begin{enumerate}
        \item 10\%
        \item 80\%
        \item 90\%
        \item 100\%
    \end{enumerate}
    
    \item ¿Cuándo es recomendable usar Leave-One-Out cross validation?
    \begin{enumerate}
        \item Con datasets muy grandes
        \item Con datasets muy pequeños
        \item Siempre que sea posible
        \item Nunca
    \end{enumerate}
    
    \item ¿Qué es stratified cross validation?
    \begin{enumerate}
        \item Una técnica para datasets grandes
        \item Una forma de mantener la proporción de clases en cada fold
        \item Un método más rápido de validación
        \item Una técnica solo para regresión
    \end{enumerate}
    
    \item ¿Cuál es el trade-off principal al elegir el número de folds?
    \begin{enumerate}
        \item Sesgo vs Varianza
        \item Velocidad vs Precisión
        \item Complejidad vs Simplicidad
        \item Todas las anteriores
    \end{enumerate}
\end{enumerate}

\subsection{Ejercicios Prácticos}

\subsubsection{Matriz de Confusión y Métricas Básicas}

Un modelo de clasificación binaria para detectar fraude en transacciones tiene los siguientes resultados en un conjunto de prueba de 1000 transacciones:

\begin{center}
\begin{tikzpicture}
\draw (0,0) rectangle (4,4);
\draw (0,2) -- (4,2);
\draw (2,0) -- (2,4);

\node[above] at (2,4) {\textbf{Real}};
\node[left] at (0,2) {\textbf{Predicho}};

\fill[green!30] (0,2) rectangle (2,4);
\fill[red!30] (2,0) rectangle (4,2);
\fill[blue!30] (0,0) rectangle (2,2);
\fill[yellow!30] (2,2) rectangle (4,4);

\node at (1,3) {45};
\node at (3,3) {15};
\node at (1,1) {5};
\node at (3,1) {935};
\end{tikzpicture}
\end{center}

\begin{enumerate}
    \item[a.] Calcular todas las métricas básicas: accuracy, precisión, recall, especificidad, F1-Score.
    \item[b.] ¿Qué puede concluir sobre el rendimiento del modelo?
    \item[c.] ¿Es la accuracy una buena métrica para este problema? Justificar.
    \item[d.] ¿Qué métrica sería más importante para este caso de uso y por qué?
\end{enumerate}

\subsubsection{Predicción de Precios}

Un modelo para predecir el precio de casas (en miles de dólares) tiene los siguientes resultados en un conjunto de prueba:

\begin{table}[H]
    \centering
    \begin{tabular}{|c|c|c|}
        \hline
        \textbf{Casa} & \textbf{Precio Real} & \textbf{Precio Predicho} \\
        \hline
        1 & 300 & 280 \\
        2 & 450 & 470 \\
        3 & 200 & 190 \\
        4 & 600 & 620 \\
        5 & 350 & 330 \\
        6 & 500 & 480 \\
        7 & 250 & 260 \\
        8 & 400 & 420 \\
        \hline
    \end{tabular}
\end{table}

\begin{enumerate}
    \item[a.] Calcular MSE, RMSE, MAE y R² para este modelo.
    \item[b.] ¿Qué métrica usarías para reportar el rendimiento a un cliente no técnico?
    \item[c.] Si una casa cuesta 0 dólares (regalo), ¿qué problema tendrías con el MAPE?
    \item[d.] ¿Qué estrategia usarías para manejar este problema?
\end{enumerate}

\subsubsection{Comparación de Modelos de Predicción}

Tenés dos modelos para predecir ventas mensuales (en miles de unidades):

\textbf{Modelo A:} MSE = 25, MAE = 3.5, R² = 0.85 \\
\textbf{Modelo B:} MSE = 30, MAE = 3.2, R² = 0.82

\begin{enumerate}
    \item[a.] ¿Cuál modelo elegirías si los errores grandes son muy costosos?
    \item[b.] ¿Cuál modelo elegirías si todos los errores tienen el mismo costo?
    \item[c.] ¿Qué información adicional necesitarías para tomar una decisión más informada?
    \item[d.] Si el R² del Modelo A fuera 0.95, ¿cambiaría tu respuesta en (a)?
\end{enumerate}

\subsubsection{Comparación de Modelos}

Tenés dos modelos para clasificar emails como spam o no spam. Los resultados son: 
\begin{itemize}
    \item \textbf{Modelo A:} VP = 180, VN = 820, FP = 30, FN = 20
    \item \textbf{Modelo B:} VP = 160, VN = 830, FP = 20, FN = 40
\end{itemize}

\begin{enumerate}
    \item[a.] Calcular precisión, recall y F1-Score para ambos modelos.
    \item[b.] ¿Cuál modelo elegirías y por qué?
    \item[c.] ¿Cambiaría tu elección si el costo de un falso negativo fuera 10 veces mayor que el de un falso positivo?
\end{enumerate}

\subsubsection{Matriz de Confusión Multiclase}

Un modelo clasifica el clima en tres categorías: Soleado (S), Nublado (N), Lluvioso (L). La matriz de confusión es:

\begin{center}
\begin{tikzpicture}
\node[above] at (3,7) {\textbf{Real}};
\node[left] at (-1,3) {\textbf{Predicho}};

\draw (0,0) rectangle (6,6);
\draw (0,2) -- (6,2);
\draw (0,4) -- (6,4);
\draw (2,0) -- (2,6);
\draw (4,0) -- (4,6);

\fill[green!30] (0,4) rectangle (2,6);
\fill[green!30] (2,2) rectangle (4,4);
\fill[green!30] (4,0) rectangle (6,2);

\node at (1,5) {85};
\node at (3,5) {10};
\node at (5,5) {5};
\node at (1,3) {8};
\node at (3,3) {82};
\node at (5,3) {10};
\node at (1,1) {7};
\node at (3,1) {12};
\node at (5,1) {81};

\node at (-0.5,5) {S};
\node at (-0.5,3) {N};
\node at (-0.5,1) {L};

\node at (1,6.5) {S};
\node at (3,6.5) {N};
\node at (5,6.5) {L};
\end{tikzpicture}
\end{center}

\begin{enumerate}
    \item[a.] Calcular la accuracy total del modelo.
    \item[b.] Calcular la precisión, recall y F1-Score para cada clase.
    \item[c.] ¿Qué clase es la más difícil de predecir para el modelo?
    \item[d.] Calcular las métricas macro y micro average.
\end{enumerate}

\subsubsection{Curva ROC}

Un modelo de clasificación binaria produce las siguientes probabilidades para 8 casos de prueba:

\begin{table}[H]
    \centering
    \begin{tabular}{|c|c|c|}
        \hline
        \textbf{Caso} & \textbf{Clase Real} & \textbf{Probabilidad Positiva} \\
        \hline
        1 & Positivo & 0.9 \\
        \hline
        2 & Negativo & 0.8 \\
        \hline
        3 & Positivo & 0.7 \\
        \hline
        4 & Positivo & 0.6 \\
        \hline
        5 & Negativo & 0.4 \\
        \hline
        6 & Negativo & 0.3 \\
        \hline
        7 & Positivo & 0.2 \\
        \hline
        8 & Negativo & 0.1 \\
        \hline
    \end{tabular}
\end{table}

\begin{enumerate}
    \item[a.] Para umbrales 0.5, 0.7 y 0.9, calcular TPR y FPR.
    \item[b.] ¿Qué umbral recomendarías si querés minimizar falsos positivos?
    \item[c.] ¿Qué umbral recomendarías si querés minimizar falsos negativos?
    \item[d.] Estimar el valor aproximado del AUC.
\end{enumerate}

\subsubsection{Problema Médico}

Un hospital desarrolla un modelo para detectar neumonía en radiografías. El dataset tiene las siguientes características:
\begin{itemize}
    \item 10,000 radiografías en total
    \item 500 casos positivos (neumonía)
    \item 9,500 casos negativos (sin neumonía)
\end{itemize}

El modelo A tiene una accuracy del 96\% pero detecta solo el 40\% de los casos de neumonía.
El modelo B tiene una accuracy del 92\% pero detecta el 85\% de los casos de neumonía.

\begin{enumerate}
    \item[a.] Calcular las matrices de confusión aproximadas para ambos modelos.
    \item[b.] ¿Cuál modelo elegirías para uso clínico? Justificar.
    \item[c.] ¿Qué estrategias podrías usar para mejorar el modelo seleccionado?
\end{enumerate}

\subsubsection{Implementación en Python}

Implementar en Python un análisis completo de validación de modelos que incluya:

\begin{lstlisting}[language=Python]
import numpy as np
import matplotlib.pyplot as plt
from sklearn.datasets import make_classification
from sklearn.model_selection import train_test_split
from sklearn.ensemble import RandomForestClassifier
from sklearn.linear_model import LogisticRegression
from sklearn.metrics import (classification_report, confusion_matrix, 
                            roc_curve, auc, precision_recall_curve)

# Generar dataset desbalanceado
X, y = make_classification(n_samples=1000, n_features=20, 
                          n_classes=2, weights=[0.9, 0.1], 
                          random_state=42)

# Completar el codigo para:
# 1. Dividir en train/test
# 2. Entrenar dos modelos diferentes
# 3. Calcular y comparar metricas
# 4. Graficar curvas ROC y PR
# 5. Crear matrices de confusion
# 6. Analizar cual modelo es mejor para este caso
\end{lstlisting}

\subsubsection{Caso de Negocio}

Una empresa de e-commerce quiere implementar un sistema de detección de reseñas falsas. Tenés los siguientes datos de rendimiento de tres modelos:

\begin{table}[H]
    \centering
    \begin{tabular}{|c|c|c|c|c|}
        \hline
        \textbf{Modelo} & \textbf{Accuracy} & \textbf{Precisión} & \textbf{Recall} & \textbf{AUC} \\
        \hline
        Naive Bayes & 0.78 & 0.65 & 0.72 & 0.76 \\
        \hline
        Random Forest & 0.82 & 0.71 & 0.68 & 0.81 \\
        \hline
        SVM & 0.80 & 0.73 & 0.64 & 0.79 \\
        \hline
    \end{tabular}
\end{table}

\begin{enumerate}
    \item[a.] ¿Qué modelo recomendarías si el objetivo principal es no bloquear reseñas legítimas?
    \item[b.] ¿Qué modelo recomendarías si el objetivo principal es mantener la calidad alta de reseñas?
    \item[c.] ¿Qué información adicional necesitarías para tomar una mejor decisión?
    \item[d.] Proponer una estrategia que combine múltiples modelos.
\end{enumerate}

\subsection{Ejercicios Prácticos (Cross Validation)}

\subsubsection{Cálculo de Modelos a Entrenar}

Considerar los siguientes escenarios e indicar cuántos árboles es necesario entrenar en cada caso:

\begin{enumerate}
    \item En el entrenamiento de un modelo CART se usa 5-fold cross validation para seleccionar el hiperparámetro max\_depth de la siguiente grilla: [none, 10, 20, 30]
    
    \item En el entrenamiento de un modelo CART se usa 5-fold cross validation para seleccionar la mejor combinación de los siguientes hiperparámetros:
    \begin{itemize}
        \item max\_depth: [None, 10, 20, 30, 40, 50]
        \item min\_samples\_split: [2, 5, 10]
        \item min\_samples\_leaf: [1, 2, 4]
    \end{itemize}
    
    \item En el entrenamiento de un modelo Random Forest se usa 5-fold cross validation para seleccionar la mejor combinación de los siguientes hiperparámetros:
    \begin{itemize}
        \item n\_estimators: [100, 200, 500, 1000]
        \item max\_depth: [10, 20, 30]
        \item min\_samples\_split: [2, 5, 10]
    \end{itemize}
\end{enumerate}

\subsubsection{Interpretación de Resultados}

Tener en cuenta los siguientes resultados de la cross validation del hiperparámetro n\_estimators para un modelo CART. ¿Qué valor del hiperparámetro elegirían? ¿Por qué?

\begin{lstlisting}[language=Python]
{
'mean_fit_time': array([7.5903286, 12.5109745, 30.5498374, 61.4766215]),
'std_fit_time': array([2.8521561, 0.3776906, 0.5143247, 0.6323429]),
'mean_score_time': array([0.1880916, 0.0871274, 0.9984307, 1.9107995]),
'std_score_time': array([0.0091651, 0.0154946, 0.1213947, 0.1533748]),
'param_n_estimators': masked_array(data=[100, 200, 500, 1000]),
'params': [{'n_estimators': 100},
           {'n_estimators': 200},
           {'n_estimators': 500},
           {'n_estimators': 1000}],
'split0_test_score': array([0.7995629, 0.7997190, 0.7998751, 0.7998751]),
'split1_test_score': array([0.7995629, 0.7997190, 0.7997190, 0.7997190]),
'split2_test_score': array([0.7995316, 0.7998438, 0.7998438, 0.7998438]),
'split3_test_score': array([0.7993754, 0.7996877, 0.7996877, 0.7998438]),
'split4_test_score': array([0.7996877, 0.7998438, 0.7998438, 0.7998438]),
'mean_test_score': array([0.7995754, 0.799762, 0.7997939, 0.7998251]),
'std_test_score': array([1.2270604e-04, 6.7250609e-05, 7.5455567e-05, 5.4430172e-05]),
'rank_test_score': array([4, 3, 2, 1], dtype=int32)
}
\end{lstlisting}

\subsubsection{KNN y Cross Validation}

Considerar el siguiente conjunto de datos donde se quiere predecir la etiqueta, + o - en función de la variable X. Con este conjunto de datos se quiere entrenar un modelo knn, para el cual el hiperparámetro k se seleccionará por cross-validation:

\begin{table}[H]
    \centering
    \begin{tabular}{|c|c|c|c|c|c|c|c|c|c|c|}
        \hline
        & \textbf{1} & \textbf{2} & \textbf{3} & \textbf{4} & \textbf{5} & \textbf{6} & \textbf{7} & \textbf{8} & \textbf{9} & \textbf{10} \\
        \hline
        \textbf{X} & -0.1 & 0.7 & 1.0 & 1.6 & 2.0 & 2.5 & 3.0 & 3.5 & 4.1 & 4.9 \\
        \hline
        \textbf{Y} & - & + & + & - & + & + & - & - & + & + \\
        \hline
    \end{tabular}
\end{table}

\begin{enumerate}
    \item Dar el accuracy de cross-validation por leave-one-out (es decir 10-folds cross-validation) de 1 vecino más cercano para la muestra de entrenamiento.
    \item Dar el accuracy de cross-validation por leave-one-out (es decir 10-folds cross-validation) de 3 vecinos más cercanos para la muestra de entrenamiento.
    \item ¿Qué valor de k seleccionas?
\end{enumerate}

\subsubsection{Implementación de K-Fold}

Implementá k-fold cross validation desde cero en Python para evaluar un modelo de regresión lineal:

\begin{lstlisting}[language=Python]
import numpy as np
from sklearn.linear_model import LinearRegression
from sklearn.metrics import mean_squared_error

def k_fold_cv(X, y, k=5, model_class=LinearRegression):
    """
    Implementa k-fold cross validation
    
    Parameters:
    X: features
    y: target
    k: numero de folds
    model_class: clase del modelo a evaluar
    
    Returns:
    scores: lista de scores para cada fold
    """
    # Tu codigo aca
    pass

# Datos de ejemplo
X = np.random.randn(100, 3)
y = X[:, 0] + 2*X[:, 1] - X[:, 2] + np.random.randn(100)*0.1

# Evaluar modelo
scores = k_fold_cv(X, y, k=5)
print(f"MSE promedio: {np.mean(scores):.4f}")
print(f"Desviacion estandar: {np.std(scores):.4f}")
\end{lstlisting}

\subsubsection{Stratified Cross Validation}

Tenemos un dataset desbalanceado con las siguientes clases:

\begin{table}[H]
    \centering
    \begin{tabular}{cc}
        \toprule
        \textbf{Clase} & \textbf{Cantidad} \\
        \midrule
        A & 800 \\
        B & 150 \\
        C & 50 \\
        \bottomrule
    \end{tabular}
\end{table}

\begin{enumerate}
    \item ¿Por qué es importante usar stratified cross validation en este caso?
    \item ¿Cuántas observaciones de cada clase habría en cada fold si usamos 5-fold stratified CV?
    \item ¿Qué pasaría si usáramos k-fold simple sin estratificación?
\end{enumerate}

\subsubsection{Grid Search con Cross Validation}

Implementá grid search con cross validation para optimizar hiperparámetros de un SVM:

\begin{lstlisting}[language=Python]
from sklearn.svm import SVC
from sklearn.model_selection import GridSearchCV
from sklearn.datasets import make_classification

# Generar datos
X, y = make_classification(n_samples=1000, n_features=10, 
                          n_classes=2, random_state=42)

# Definir grilla de parametros
param_grid = {
    'C': [0.1, 1, 10, 100],
    'gamma': ['scale', 'auto', 0.001, 0.01, 0.1],
    'kernel': ['rbf', 'linear']
}

# Tu codigo aca para implementar grid search con CV
\end{lstlisting}

\subsubsection{Comparación de Modelos}

Usando cross validation, compará el rendimiento de los siguientes modelos en un problema de clasificación:

\begin{enumerate}
    \item Logistic Regression
    \item Random Forest
    \item SVM
    \item KNN
\end{enumerate}

¿Qué métrica usarías? ¿Cómo determinarías si las diferencias son estadísticamente significativas?

\subsubsection{Análisis de Varianza}

\begin{enumerate}
    \item ¿Por qué es importante reportar tanto la media como la desviación estándar de los scores de CV?
    \item ¿Qué indica una alta varianza en los scores de cross validation?
    \item ¿Cómo podríamos reducir la varianza en los resultados de CV?
\end{enumerate}

\newpage

\section{Soluciones}

\subsection{Ejercicios Conceptuales}

\subsubsection{Verdadero o Falso}

\begin{enumerate}
    \item \textbf{Falso.} Accuracy puede ser engañosa con clases desbalanceadas.
    \item \textbf{Verdadero.} AUC < 0.5 es peor que aleatorio (0.5).
    \item \textbf{Verdadero.} Precisión = VP / (VP + FP).
    \item \textbf{Falso.} Recall + Especificidad no suman 1. Especificidad + FPR suman 1.
    \item \textbf{Verdadero.} Es más importante no perder casos positivos (enfermos).
    \item \textbf{Falso.} ROC grafica TPR vs FPR. PR grafica Precisión vs Recall.
    \item \textbf{Verdadero.} F1 es la media armónica de precisión y recall.
    \item \textbf{Falso.} También se usa para problemas multiclase.
    \item \textbf{Falso.} AUC = 0.5 indica clasificador aleatorio, no perfecto.
    \item \textbf{Verdadero.} PR es más sensible a clases minoritarias.
    \item \textbf{Falso.} Depende del contexto y costos del problema.
    \item \textbf{Verdadero.} Se prefiere no condenar inocentes.
    \item \textbf{Verdadero.} FDR = 1 - Precisión.
    \item \textbf{Falso.} Hay múltiples tipos de errores, pero el concepto se extiende.
    \item \textbf{Verdadero.} Especialmente con clases desbalanceadas.
    \item \textbf{Verdadero.} MSE eleva al cuadrado los errores, penalizando más los grandes.
    \item \textbf{Verdadero.} R² negativo indica que el modelo es peor que predecir la media.
    \item \textbf{Verdadero.} RMSE = $\sqrt{MSE}$, mantiene las unidades originales.
    \item \textbf{Verdadero.} MAPE tiene división por cero cuando $y_i = 0$.
    \item \textbf{Verdadero.} R² mide la proporción de varianza explicada.
\end{enumerate}

\subsubsection{Verdadero o Falso (Cross Validation)}

\begin{enumerate}
    \item \textbf{Verdadero.} CV proporciona múltiples estimaciones del rendimiento, reduciendo la dependencia de una única división train/test.
    \item \textbf{Falso.} Los folds pueden tener tamaños ligeramente diferentes si n no es divisible por k.
    \item \textbf{Falso.} LOOCV puede tener alta varianza y ser computacionalmente costoso.
    \item \textbf{Verdadero.} CV se usa tanto para model selection como para model evaluation.
    \item \textbf{Verdadero.} En 5-fold CV, cada observación se usa 4 veces para entrenar y 1 vez para validar.
    \item \textbf{Verdadero.} Mantiene la proporción de clases en cada fold.
    \item \textbf{Falso.} CV ayuda a detectar overfitting, pero no lo reduce directamente.
    \item \textbf{Falso.} Más folds aumenta la varianza y el costo computacional.
\end{enumerate}

\subsubsection{Multiple Choice}

\begin{enumerate}
    \item \textbf{b)} Precisión se enfoca en predicciones positivas correctas, recall en casos positivos reales detectados.
    \item \textbf{c)} Recall - es crucial no perder casos de cáncer.
    \item \textbf{b)} Un umbral de clasificación diferente.
    \item \textbf{a)} Accuracy - un modelo que predice todo negativo tendría 95\% accuracy.
    \item \textbf{d)} Cuando las clases están desbalanceadas.
    \item \textbf{b)} 80\% de probabilidad de ranquear correctamente un par positivo-negativo.
    \item \textbf{a)} MSE - eleva al cuadrado los errores, amplificando el impacto de valores atípicos.
    \item \textbf{a)} El modelo explica el 30\% de la varianza de la variable objetivo.
    \item \textbf{b)} Es fácil de interpretar como porcentaje, aunque tiene problemas con ceros.
\end{enumerate}

\subsubsection{Multiple Choice (Cross Validation)}

\begin{enumerate}
    \item \textbf{b)} Usa toda la data tanto para entrenamiento como para validación en diferentes momentos.
    \item \textbf{c)} 90\% (9 de los 10 folds se usan para entrenamiento).
    \item \textbf{b)} Con datasets muy pequeños, cuando queremos usar todos los datos posibles.
    \item \textbf{b)} Mantiene la proporción de clases en cada fold.
    \item \textbf{d)} Todas las anteriores: más folds aumentan sesgo pero reducen varianza, son más lentos pero más precisos, y más complejos pero potencialmente mejores.
\end{enumerate}

\subsection{Ejercicios Prácticos}

\subsubsection{Matriz de Confusión y Métricas Básicas}

\textbf{a)} Métricas calculadas:
\begin{align*}
\text{Accuracy} &= \frac{45 + 935}{1000} = 0.98 \\
\text{Precisión} &= \frac{45}{45 + 15} = 0.75 \\
\text{Recall} &= \frac{45}{45 + 5} = 0.90 \\
\text{Especificidad} &= \frac{935}{935 + 15} = 0.984 \\
\text{F1} &= 2 \times \frac{0.75 \times 0.90}{0.75 + 0.90} = 0.818
\end{align*}

\textbf{b)} El modelo tiene alta accuracy y recall, pero precisión moderada.

\textbf{c)} No, accuracy puede ser engañosa porque hay pocas transacciones fraudulentas.

\textbf{d)} Recall es más importante para no perder casos de fraude.

\subsubsection{Predicción de Precios}

\textbf{a)} Cálculos:
\begin{align*}
\text{MSE} &= \frac{1}{8}[(300-280)^2 + (450-470)^2 + (200-190)^2 + (600-620)^2 + (350-330)^2\\ &+ (500-480)^2 + (250-260)^2 + (400-420)^2] \\
&= \frac{1}{8}[400 + 400 + 100 + 400 + 400 + 400 + 100 + 400] = 325 \\
\text{RMSE} &= \sqrt{325} = 18.03 \\
\text{MAE} &= \frac{1}{8}[20 + 20 + 10 + 20 + 20 + 20 + 10 + 20] = 17.5 \\
\text{R²} &= 1 - \frac{325}{varianza\_real} = 1 - \frac{325}{6500} = 0.95 \\
\text{Donde: } \bar{y} &= \frac{300+450+200+600+350+500+250+400}{8} = 381.25 \\
\text{varianza\_real} &= (300-381.25)^2 + (450-381.25)^2 + ... + (400-381.25)^2 = 6500
\end{align*}

\textbf{b)} RMSE o MAE - son fáciles de interpretar en las mismas unidades.

\textbf{c)} División por cero en el denominador del MAPE.

\textbf{d)} Reemplazar ceros por valor muy pequeño, filtrar ceros, o usar sMAPE.

\subsubsection{Comparación de Modelos de Predicción}

\textbf{a)} Modelo A - tiene menor MSE (25 vs 30), penaliza menos los errores grandes.

\textbf{b)} Modelo B - tiene menor MAE (3.2 vs 3.5), trata todos los errores por igual.

\textbf{c)} Costos específicos de errores, distribución de los datos, interpretabilidad.

\textbf{d)} Sí, R² = 0.95 indicaría que el Modelo A explica mucha más varianza, compensando el mayor MAE.

\subsubsection{Comparación de Modelos}

\textbf{Modelo A:}
\begin{align*}
\text{Precisión} &= \frac{180}{210} = 0.857 \\
\text{Recall} &= \frac{180}{200} = 0.90 \\
\text{F1} &= 0.878
\end{align*}

\textbf{Modelo B:}
\begin{align*}
\text{Precisión} &= \frac{160}{180} = 0.889 \\
\text{Recall} &= \frac{160}{200} = 0.80 \\
\text{F1} &= 0.842
\end{align*}

Modelo A tiene mejor recall y F1. Si los falsos negativos cuestan 10 veces más, definitivamente elegir Modelo A.

\subsubsection{Matriz de Confusión Multiclase}

\textbf{a.} Accuracy = (85 + 82 + 81) / 300 = 0.827

\textbf{b.} Calculando por clase:
\begin{itemize}
    \item Soleado: Precisión = 0.85, Recall = 0.85, F1 = 0.85
    \item Nublado: Precisión = 0.82, Recall = 0.782, F1 = 0.80
    \item Lluvioso: Precisión = 0.844, Recall = 0.81, F1 = 0.827
\end{itemize}

\textbf{c.} Nublado es la más difícil (menor F1).

\subsubsection{Curva ROC}

\textbf{a.} Para cada umbral:
\begin{itemize}
    \item Umbral 0.5: TPR = 0.75, FPR = 0.25
    \item Umbral 0.7: TPR = 0.75, FPR = 0.25
    \item Umbral 0.9: TPR = 0.25, FPR = 0.25
\end{itemize}

\textbf{b.} Para minimizar FP: umbral 0.9 \\
\textbf{c.} Para minimizar FN: umbral bajo (0.1) \\
\textbf{d.} AUC aproximado: 0.75

\subsubsection{Problema Médico}

\textbf{a.} Modelo A: VP=200, FN=300, FP=100, VN=9400. Modelo B: VP=425, FN=75, FP=725, VN=8775 \\
\textbf{b.} Modelo B porque detecta más casos de neumonía, crucial en medicina. \\
\textbf{c.} Estrategias: datos más balanceados, ensemble, ajuste de umbral.

\subsubsection{Implementación en Python}

\begin{lstlisting}[language=Python]
import numpy as np
import pandas as pd
from sklearn.datasets import make_classification
from sklearn.model_selection import train_test_split
from sklearn.linear_model import LogisticRegression
from sklearn.ensemble import RandomForestClassifier
from sklearn.metrics import classification_report, confusion_matrix, roc_curve, auc, precision_recall_curve
import matplotlib.pyplot as plt

# Generar un dataset desbalanceado
X, y = make_classification(n_samples=5000, n_features=20, n_informative=2, n_redundant=10,
                           n_clusters_per_class=1, weights=[0.9, 0.1], flip_y=0, random_state=42)

# Division train/test
X_train, X_test, y_train, y_test = train_test_split(X, y, test_size=0.3, stratify=y, random_state=42)

# Entrenamiento de modelos
lr = LogisticRegression()
rf = RandomForestClassifier()

lr.fit(X_train, y_train)
rf.fit(X_train, y_train)

# Predicciones
y_pred_lr = lr.predict(X_test)
y_pred_rf = rf.predict(X_test)
y_proba_lr = lr.predict_proba(X_test)[:,1]
y_proba_rf = rf.predict_proba(X_test)[:,1]

# Calculo de metricas
print("Logistic Regression")
print(confusion_matrix(y_test, y_pred_lr))
print(classification_report(y_test, y_pred_lr, digits=3))

print("Random Forest")
print(confusion_matrix(y_test, y_pred_rf))
print(classification_report(y_test, y_pred_rf, digits=3))

# Curva ROC
fpr_lr, tpr_lr, _ = roc_curve(y_test, y_proba_lr)
fpr_rf, tpr_rf, _ = roc_curve(y_test, y_proba_rf)
auc_lr = auc(fpr_lr, tpr_lr)
auc_rf = auc(fpr_rf, tpr_rf)

plt.figure()
plt.plot(fpr_lr, tpr_lr, label="Logistic Reg (AUC = %.2f)" % auc_lr)
plt.plot(fpr_rf, tpr_rf, label="Random Forest (AUC = %.2f)" % auc_rf)
plt.plot([0,1],[0,1],'k--')
plt.xlabel("FPR")
plt.ylabel("TPR")
plt.title("Curva ROC")
plt.legend()
plt.show()

# Curva Precision-Recall
prec_lr, rec_lr, _ = precision_recall_curve(y_test, y_proba_lr)
prec_rf, rec_rf, _ = precision_recall_curve(y_test, y_proba_rf)

plt.figure()
plt.plot(rec_lr, prec_lr, label="Logistic Reg")
plt.plot(rec_rf, prec_rf, label="Random Forest")
plt.xlabel("Recall")
plt.ylabel("Precision")
plt.title("Curva Precision-Recall")
plt.legend()
plt.show()

# Analisis comparativo
# En este caso, mirar F1, recall y precision para la clase minoritaria (1)
# Tambien comparar AUC y la curva PR, que es mas informativa en datasets desbalanceados
\end{lstlisting}

\subsubsection{Caso de Negocio}

\textbf{a.} Naive Bayes (mayor recall para no bloquear reseñas legítimas). \\
\textbf{b.} SVM (mayor precisión para mantener calidad). \\
\textbf{c.} Costos de errores, volumen de reseñas, recursos disponibles. \\
\textbf{d.} Ensemble que combine todo con pesos según objetivos del negocio.

\subsection{Ejercicios Prácticos (Cross Validation)}

\subsubsection{Cálculo de Modelos a Entrenar}

\begin{enumerate}
    \item \textbf{CART con 5-fold CV y 4 valores de max\_depth:}
    \begin{itemize}
        \item 4 valores de hiperparámetro × 5 folds = \textbf{20 árboles}
    \end{itemize}
    
    \item \textbf{CART con 5-fold CV y múltiples hiperparámetros:}
    \begin{itemize}
        \item max\_depth: 6 valores
        \item min\_samples\_split: 3 valores  
        \item min\_samples\_leaf: 3 valores
        \item Total combinaciones: 6 × 3 × 3 = 54
        \item 54 combinaciones × 5 folds = \textbf{270 árboles}
    \end{itemize}
    
    \item \textbf{Random Forest con 5-fold CV:}
    \begin{itemize}
        \item n\_estimators: 4 valores
        \item max\_depth: 3 valores
        \item min\_samples\_split: 3 valores
        \item Total combinaciones: 4 × 3 × 3 = 36
        \item 36 combinaciones × 5 folds = 180 Random Forests
        \item Cada Random Forest tiene su propio n\_estimators (100, 200, 500, o 1000)
        \item Total de árboles individuales = 180 × promedio de n\_estimators
        \item Promedio de n\_estimators = (100+200+500+1000)/4 = 450
        \item \textbf{Total aproximado: 81,000 árboles individuales}
    \end{itemize}
\end{enumerate}

\subsubsection{Interpretación de Resultados}

Analizando los resultados:

\begin{itemize}
    \item \textbf{n\_estimators = 100:} score = 0.7996, tiempo = 7.59s, std = 1.23e-04
    \item \textbf{n\_estimators = 200:} score = 0.7998, tiempo = 12.51s, std = 6.73e-05
    \item \textbf{n\_estimators = 500:} score = 0.7998, tiempo = 30.55s, std = 7.55e-05
    \item \textbf{n\_estimators = 1000:} score = 0.7999, tiempo = 61.48s, std = 5.44e-05
\end{itemize}

\textbf{Recomendación: n\_estimators = 100}

\vspace{1em}

\textbf{Justificación:}
\begin{itemize}
    \item La diferencia de score entre 100 y 200 es solo 0.0002 (0.7996 vs 0.7998)
    \item Esta diferencia es muy pequeña y probablemente no sea estadísticamente significativa
    \item El tiempo aumenta 65\% (de 7.59s a 12.51s) para una ganancia mínima
    \item La desviación estándar es ligeramente mayor en 100, pero aún muy baja
    \item En términos prácticos, 0.0002 de diferencia no justifica el costo computacional adicional
\end{itemize}

\subsubsection{KNN y Cross Validation}

\textbf{Dataset ordenado por X:}
\begin{table}[H]
    \centering
    \begin{tabular}{|c|c|c|c|c|c|c|c|c|c|c|}
        \hline
        \textbf{Punto} & 1 & 2 & 3 & 4 & 5 & 6 & 7 & 8 & 9 & 10 \\
        \hline
        \textbf{X} & -0.1 & 0.7 & 1.0 & 1.6 & 2.0 & 2.5 & 3.0 & 3.5 & 4.1 & 4.9 \\
        \hline
        \textbf{Y} & - & + & + & - & + & + & - & - & + & + \\
        \hline
    \end{tabular}
\end{table}

\textbf{a) LOOCV con k=1 (1-NN):}

Para cada punto, encontramos su vecino más cercano y vemos si coincide la clase:

\begin{itemize}
    \item Punto 1 (X=-0.1, Y=-): vecino más cercano es punto 2 (Y=+) → \textbf{Error}
    \item Punto 2 (X=0.7, Y=+): vecino más cercano es punto 3 (Y=+) → \textbf{Correcto}  
    \item Punto 3 (X=1.0, Y=+): vecino más cercano es punto 2 (Y=+) → \textbf{Correcto}
    \item Punto 4 (X=1.6, Y=-): vecino más cercano es punto 5 (Y=+) → \textbf{Error}
    \item Punto 5 (X=2.0, Y=+): vecino más cercano es punto 4 (Y=-) → \textbf{Error}
    \item Punto 6 (X=2.5, Y=+): vecino más cercano es punto 5 (Y=+) → \textbf{Correcto}
    \item Punto 7 (X=3.0, Y=-): vecino más cercano es punto 6 (Y=+) → \textbf{Error}
    \item Punto 8 (X=3.5, Y=-): vecino más cercano es punto 7 (Y=-) → \textbf{Correcto}
    \item Punto 9 (X=4.1, Y=+): vecino más cercano es punto 8 (Y=-) → \textbf{Error}
    \item Punto 10 (X=4.9, Y=+): vecino más cercano es punto 9 (Y=+) → \textbf{Correcto}
\end{itemize}

\textbf{Accuracy k=1: 6/10 = 0.60}

\vspace{1em}

\textbf{b) LOOCV con k=3 (3-NN):}

Para cada punto, encontramos sus 3 vecinos más cercanos y usamos votación por mayoría:

\begin{itemize}
    \item Punto 1: vecinos [2,3,4] → clases [+,+,-] → mayoría + → \textbf{Error}
    \item Punto 2: vecinos [1,3,4] → clases [-,+,-] → mayoría - → \textbf{Error}
    \item Punto 3: vecinos [2,4,5] → clases [+,-,+] → mayoría + → \textbf{Correcto}
    \item Punto 4: vecinos [3,5,6] → clases [+,+,+] → mayoría + → \textbf{Error}
    \item Punto 5: vecinos [4,6,7] → clases [-,+,-] → mayoría - → \textbf{Error}
    \item Punto 6: vecinos [5,7,8] → clases [+,-,-] → mayoría - → \textbf{Error}
    \item Punto 7: vecinos [6,8,9] → clases [+,-,+] → mayoría + → \textbf{Error}
    \item Punto 8: vecinos [7,9,10] → clases [-,+,+] → mayoría + → \textbf{Error}
    \item Punto 9: vecinos [8,10,7] → clases [-,+,-] → mayoría - → \textbf{Error}
    \item Punto 10: vecinos [9,8,7] → clases [+,-,-] → mayoría - → \textbf{Error}
\end{itemize}

\textbf{Accuracy k=3: 1/10 = 0.10}

\vspace{1em}

\textbf{c) Selección de k:}
Elegiría \textbf{k=1} porque tiene mejor accuracy (0.60 vs 0.10). Con este dataset pequeño y ruidoso, k=3 incluye demasiados vecinos que añaden ruido a la predicción.

\subsubsection{Implementación de K-Fold}

\begin{lstlisting}[language=Python]
import numpy as np
from sklearn.linear_model import LinearRegression
from sklearn.metrics import mean_squared_error

def k_fold_cv(X, y, k=5, model_class=LinearRegression):
    """
    Implementa k-fold cross validation
    """
    n_samples = len(X)
    fold_size = n_samples // k
    indices = np.arange(n_samples)
    np.random.shuffle(indices)  # mezclar indices
    
    scores = []
    
    for i in range(k):
        # Definir indices de validacion
        start_idx = i * fold_size
        end_idx = start_idx + fold_size if i < k-1 else n_samples
        val_indices = indices[start_idx:end_idx]
        train_indices = np.concatenate([indices[:start_idx], 
                                       indices[end_idx:]])
        
        # Dividir datos
        X_train, X_val = X[train_indices], X[val_indices]
        y_train, y_val = y[train_indices], y[val_indices]
        
        # Entrenar modelo
        model = model_class()
        model.fit(X_train, y_train)
        
        # Predecir y evaluar
        y_pred = model.predict(X_val)
        score = mean_squared_error(y_val, y_pred)
        scores.append(score)
    
    return scores
\end{lstlisting}

\subsubsection{Stratified Cross Validation}

\begin{enumerate}
    \item \textbf{Importancia de stratified CV:} Con clases desbalanceadas, k-fold simple podría crear folds sin representación de clases minoritarias. Stratified CV garantiza que cada fold mantenga la proporción original de clases.
    
    \item \textbf{Observaciones por fold en 5-fold stratified CV:}
    \begin{itemize}
        \item Total: 1000 observaciones
        \item Proporción: A=80\%, B=15\%, C=5\%
        \item Por fold (200 obs): A=160, B=30, C=10
    \end{itemize}
    
    \item \textbf{Problema sin estratificación:} Algunos folds podrían tener muy pocas o ninguna observación de clases B y C, dando estimaciones sesgadas del rendimiento.
\end{enumerate}

\subsubsection{Grid Search con Cross Validation}

\begin{lstlisting}[language=Python]
from sklearn.svm import SVC
from sklearn.model_selection import GridSearchCV
from sklearn.datasets import make_classification

# Generar datos
X, y = make_classification(n_samples=1000, n_features=10, 
                          n_classes=2, random_state=42)

# Definir grilla de parametros
param_grid = {
    'C': [0.1, 1, 10, 100],
    'gamma': ['scale', 'auto', 0.001, 0.01, 0.1],
    'kernel': ['rbf', 'linear']
}

# Grid search con CV
grid_search = GridSearchCV(
    SVC(), 
    param_grid, 
    cv=5, 
    scoring='accuracy',
    n_jobs=-1,
    verbose=1
)

grid_search.fit(X, y)

print(f"Mejores parametros: {grid_search.best_params_}")
print(f"Mejor score: {grid_search.best_score_:.4f}")
\end{lstlisting}

\subsubsection{Comparación de Modelos}

\begin{lstlisting}[language=Python]
from sklearn.model_selection import cross_val_score
from sklearn.linear_model import LogisticRegression
from sklearn.ensemble import RandomForestClassifier
from sklearn.svm import SVC
from sklearn.neighbors import KNeighborsClassifier
import numpy as np

# Modelos a comparar
models = {
    'Logistic Regression': LogisticRegression(),
    'Random Forest': RandomForestClassifier(),
    'SVM': SVC(),
    'KNN': KNeighborsClassifier()
}

# Comparar modelos
results = {}
for name, model in models.items():
    scores = cross_val_score(model, X, y, cv=5, scoring='f1_weighted')
    results[name] = scores
    print(f"{name}: {scores.mean():.4f} (+/- {scores.std()*2:.4f})")

# Test estadistico (t-test pareado) para comparar modelos
from scipy import stats
model1_scores = results['Random Forest']
model2_scores = results['SVM'] 
t_stat, p_value = stats.ttest_rel(model1_scores, model2_scores)
print(f"p-value: {p_value:.4f}")
\end{lstlisting}

Usaría \textbf{F1-score weighted} para datasets desbalanceados o \textbf{accuracy} para balanceados. Para significancia estadística, usaría t-test pareado entre scores de CV.

\subsubsection{Análisis de Varianza}

\begin{enumerate}
    \item \textbf{Importancia de media y desviación:} La media indica el rendimiento esperado, la desviación indica la estabilidad del modelo. Ambas son cruciales para tomar decisiones informadas.
    
    \item \textbf{Alta varianza indica:}
    \begin{itemize}
        \item Modelo sensible a cambios en datos de entrenamiento
        \item Posible overfitting
        \item Dataset pequeño o muy heterogéneo
        \item Necesidad de regularización
    \end{itemize}
    
    \item \textbf{Reducir varianza:}
    \begin{itemize}
        \item Aumentar el dataset
        \item Usar modelos más simples
        \item Aplicar regularización
        \item Aumentar el número de folds (trade-off con costo computacional)
        \item Usar stratified CV para datasets desbalanceados
    \end{itemize}
\end{enumerate}

\end{document}
